% Fourth round: the port to Bend 2 and what its parallelism delivered.
\section{A port to Bend 2, and what its parallelism delivered}
\label{sec:bend}

Bend 2 (version 2.0.10) is a pure, total, affine functional language that compiles
through C to native code and promises parallelism for free: a parallel let
\texttt{a b = f(x) g(y)} forks, the type checker guarantees that the two calls
share nothing, and the result cannot depend on the number of threads. The
algorithm of \texttt{fastunknot} was ported to it (directory \texttt{bend/}) to
find out what that promise is worth for this problem. Nothing in this section
changes the complexity discussion: the scan is exponential in the boundary
width, in every language.

\subsection{Setting}
Bend runs on Linux, macOS and WSL only, so all Bend work ran over SSH on an
Ubuntu virtual machine with 6 virtual CPUs of an i7-3770. That CPU has no AVX2
and the distributed binary died with an illegal instruction; Bend was run from
source under the baseline build of its JavaScript runtime instead. All times
below are from that host, and the Python package was copied there for every
comparison, so the two sides of a ratio always share the hardware. Timings are
interleaved with an A/A control as in Section~\ref{sec:round3}; the control's
ratio stayed within $0.96$ to $1.05$.

\subsection{What was ported, and how it was checked}
\begin{itemize}
\item the modular Alexander filter (Wirtinger matrix, determinant modulo 65521,
  because \texttt{U32} is the only machine integer);
\item the planar geometry, the dotted cobordism algebra over $\mathbb{F}_2$ and
  Bar-Natan's scan with delooping and Gaussian elimination, that is
  \texttt{scan\_fast.py} without its caches;
\item a greedy scan order from eight starting crossings;
\item a standalone recognizer that chains the three.
\end{itemize}
Not ported: diagram simplification, the Jones filter, factorization, and the
interning of matchings with plan caches. Each layer was compared with its Python
counterpart before the next was written: 200 Alexander minors, 1920 gluings and
circle numberings, 2736 compositions and transfers, the Khovanov ranks of 160
random closures and of the timing inputs (up to rank 5898), and the verdicts of
the recognizer on 60 diagrams with up to 40 crossings, 40 of them unknots, among
the knots Conway's and the Kinoshita--Terasaka knot: no mismatch anywhere.
Sixteen laws about the control logic (for instance: of two scan orders with
equal profile the left one wins, so the work does not depend on scheduling) are
stated in \texttt{LAWS.bend} and proved in \texttt{PROOF.bend}; they were drafted
by the AI as proposals, and they do not say that a rank is right, which Bend
cannot express about \texttt{U32} arithmetic.

\subsection{What the fork delivers}
Measured on kernels first: a deep, balanced, uniform recursion reaches 5.0 on 6
threads; 32 coarse leaves reach 2.1; every switch from sequential to parallel
execution costs about 80 microseconds; forks inside the sparse elimination of
the Alexander determinant made it 8 to 15 times \emph{slower}; one fork over
four independent determinants gained nothing. Tasks are dealt once and never
moved, and there is no cancellation, so the race of scan orders, the best
parallel gain of the Rust port ($0.65$ of the time in total), cannot be written.

In the scan, the transfer of all objects of a stage through a crossing is a
map that shares nothing, and it was written as a fork over the tree of the
array. Its first effect was negative: with six threads the scan took 1.28 to
1.85 s instead of 1.04 s. The cause was found in the runtime source and the
emitted C, and then confirmed by measurement: generated code makes a non-tail
call either by pushing a stack frame or, in ``parallel mode'', by allocating
the continuation as a task on the heap; a program runs in parallel mode
wherever a fork is reachable, which for a program with one fork is
everywhere. With the forking transfer \emph{switched off} and six threads the
scan took 1.85 s: a factor 1.8 for the mere presence of a fork. The remedy is
to make every long sequential phase a task of its own, a fork next to a
trivial second child, since a task whose function reaches no fork runs in the
fast mode. Three natural ways of writing that child were silently compiled
away (inlined, constant-folded, dropped as unused), leaving no fork; it has to
be recursive, depend on run-time data, and its result has to be used. After
that: 1.02 s on one thread, 0.81 s on six; the transfer alone became 2.7 times
faster.

\subsection{Where the time goes}
\begin{center}
\begin{tabular}{lrrrr}
\toprule
 & \multicolumn{2}{c}{braid closure, rank 218} & \multicolumn{2}{c}{stress closure, rank 5898}\\
version & 1 thread & 6 threads & 1 thread & 6 threads\\
\midrule
first working scan, first-found pivots & 1.04 s & 1.28--1.85 s & 63 s & 58 s\\
sequential phases as tasks & 1.05 s & 0.89 s & 63 s & 58 s\\
least-fill pivots in batches, work lists & 0.65 s & 0.50 s & 7.0 s & 4.9 s\\
Python package, same host & 0.134 s & & 0.542 s & \\
\bottomrule
\end{tabular}
\end{center}
The largest gain was algorithmic and sequential: cancelling only pivots whose
fill-in is within a limit, and raising the limit to the least fill-in that was
skipped, brought the stress closure from 58 to 5 s. A sampling profile of the
emitted C then showed 43\% of the time in reference counting: a value used more
than once is shared and counted, and every step through a shared list bumps and
drops counts. Matchings are shared by design, and the linear lookups in them are
most of those steps. Python avoids the work altogether by interning matchings
as small integers and caching every plan (2817 plans serve 251\,406 transferred
entries on the stress closure); in Bend a cache is a value every call must
thread, the only map has string keys, and an array cannot be shared between the
branches of a fork. That layer was not ported. It would make the port faster
and, since it removes work from the one part that forks, less parallel.

As a process the native recognizer beats the Python command line where
Python's 28 ms start-up dominates (2 ms against 28 ms on an 8-crossing unknot)
and loses from about 25 crossings on (80 against 45 ms at 30 crossings);
in-process, Python is 2 to 4 times faster on the same steps.

\subsection{Verdict}
The promise holds in the narrow sense and not in the useful one. Writing a fork
is trivial and safe. But the scan alternates a parallel map with a Gaussian
elimination in which each cancellation changes the rows the next one reads; the
map is 20 to 45\% of the time, so six threads give 1.2 to 1.4 overall, on top of
a single-thread speed 4 to 13 times below the tuned Python package. The cost
model (a fork anywhere taxes all sequential code by 1.8 unless it is wrapped as
a task) is invisible in the language, and the compiler removes the wrapper
unless it is written just so. For this algorithm the measured ranking stays:
Rust, then Python, then Bend.
